% Original manuscript: Myeongjin Kang, 2026. CC BY 4.0.
\documentclass[11pt]{article}
\usepackage[T1]{fontenc}
\usepackage{lmodern,amsmath,amssymb,amsthm,mathtools}
\usepackage[margin=1in]{geometry}
\usepackage{microtype}
\usepackage[colorlinks=true,linkcolor=blue,citecolor=blue,urlcolor=blue]{hyperref}
\usepackage{enumitem}
\setlist{itemsep=3pt,topsep=5pt}
\emergencystretch=1em
\newtheorem{theorem}{Theorem}[section]
\newtheorem{proposition}[theorem]{Proposition}
\newtheorem{lemma}[theorem]{Lemma}
\newtheorem{corollary}[theorem]{Corollary}
\theoremstyle{definition}\newtheorem{definition}[theorem]{Definition}
\theoremstyle{remark}\newtheorem{remark}[theorem]{Remark}
\DeclareMathOperator{\Tr}{Tr}\DeclareMathOperator{\tr}{tr}
\DeclareMathOperator{\Vol}{Vol}\DeclareMathOperator{\rank}{rank}
\DeclareMathOperator{\Ad}{Ad}\DeclareMathOperator{\alg}{alg}
\DeclareMathOperator{\Rea}{Re}\DeclareMathOperator{\Ima}{Im}
\newcommand{\C}{\mathbb C}\newcommand{\R}{\mathbb R}
\newcommand{\N}{\mathbb N}\newcommand{\sa}{\mathrm{sa}}
\newcommand{\eps}{\varepsilon}\newcommand{\id}{\mathrm{id}}
\newcommand{\HS}{\mathrm{HS}}\newcommand{\op}{\mathrm{op}}
\newcommand{\gap}{\mathcal G}\newcommand{\Enc}{E_{\eps}}
\newcommand{\norm}[1]{\left\|#1\right\|}
\newcommand{\pin}{\texttt{fd4aeeb2ee4fc729c18d98444fed42fd0529eeeb}}
\title{A finite free-entropy gap for two self-adjoint variables}
\author{Myeongjin Kang}
\date{October 18, 2026}
\begin{document}
\maketitle
\begin{abstract}
Starting from the stipulated finite-entropy separation in the pinned OpenAI
Math preprint, we construct exactly two bounded self-adjoint elements with
$-\infty<\chi<\chi^*<\infty$. The microstates entropy is the original
operator-norm-cutoff entropy with an ordinary limsup over all matrix sizes.
For the source tuple $X$ of length $n=2^{64}+1$, put $k=2^{32}$.
A normal-slice encoding $E_\eps(X)$ satisfies
$\liminf_{\eps\downarrow0}(\chi^*(E_\eps(X))-\chi(E_\eps(X)))
\geq k^{-2}(\chi^*(X)-\chi(X))$.
The proof couples an unconditional Fisher-information estimate with a
matched microstates-volume estimate. The latter recovers matrix units
from arbitrary pair microstates after removing a vanishing fraction of
the dimension, so the upper bound includes every output size.
All analytic estimates needed for this transfer are proved below;
the source separation and explicitly named classical free-entropy
interfaces are inputs. No assertion about generation of the ambient
von Neumann algebra is made.
\end{abstract}

\section{Statement, source, and contribution}
\label{sec:intro}
All tracial von Neumann algebras in this paper have a faithful normal
tracial state. All tuples are finite and consist of bounded self-adjoint
elements. The two free entropies have the normalizations of
Section~\ref{sec:conventions}. In particular, $\chi$ never denotes a
relative entropy or an ultrafilter entropy unless explicitly decorated.

\begin{theorem}\label{thm:main}
There are a tracial von Neumann algebra $(\mathcal N,\tau_{\mathcal N})$
and two bounded self-adjoint elements $Z_1,Z_2\in\mathcal N$ such that
\[
 -\infty<\chi(Z_1,Z_2)<\chi^*(Z_1,Z_2)<\infty.
\]
Here $\chi$ uses the original operator-norm cutoff and ordinary limsup
over matrix sizes.
\end{theorem}

The source input is the main theorem and the specified construction of
\cite{source298}, at the immutable revision recorded in
Appendix~\ref{app:provenance}. Its theorems are stipulated inputs to this
work; we do not reprove or audit them. More precisely, that construction
gives
\begin{equation}\label{eq:source}
 n=2^{64}+1,\qquad X=C+\sqrt h S,\quad h>0,\qquad
 -\infty<\chi(X)\leq\chi^*(X)-\tfrac12<\infty.
\end{equation}
Here $C=(B_1,\ldots,B_p,Y)$, $p=2^{64}$, the $B_i$ are standard free
semicircular variables, and $Y$ is uniform on $p$ equally spaced points
in $[-1,1]$ in a separate tensor factor. The family $S$ is standard
semicircular and free from $C$. The source constructs bounded matrix
models. Thus $W^*(X)$ is embeddable into a tracial ultrapower of the
hyperfinite factor, and
\begin{equation}\label{eq:finite-fisher-source}
 \Phi^*(X)\leq n/h<\infty.
\end{equation}
These properties, the length $n$, and the positive finite gap are the
only information imported from the source construction.

Set $k=2^{32}$, so $n=k^2+1$. Let $D_1=\operatorname{diag}(1,\ldots,k)$
and let $D_2$ be the adjacency matrix of the path on $\{1,\ldots,k\}$.
On Hermitian pairs use $\langle A,B\rangle=\sum_{s=1}^2\tr_k(A_sB_s)$.
Let $V_1,\ldots,V_n$ be any real orthonormal basis of the orthogonal
complement of
\[
 \{(i[H,D_1],i[H,D_2]):H\in M_k(\C)^{\sa}\}.
\]
Such a complement has dimension $k^2+1$, as proved below. Define
\begin{equation}\label{eq:encoding}
 E_\eps(X)_s=D_s\otimes1+\eps\sum_{a=1}^n V_{as}\otimes X_a,
 \qquad s=1,2.
\end{equation}
The ambient algebra is $M_k(W^*(X))$ with trace $\tr_k\otimes\tau$.

\begin{theorem}[Transfer for the specified source]\label{thm:transfer}
For $X$ in \eqref{eq:source}, both entropies of $E_\eps(X)$ are finite
for every sufficiently small $\eps>0$, and
\begin{equation}\label{eq:transfer}
 \liminf_{\eps\downarrow0}\gap(E_\eps(X))\geq k^{-2}\gap(X),
 \qquad \gap(A)=\chi^*(A)-\chi(A).
\end{equation}
Consequently every sufficiently small positive encoding parameter gives
a strictly positive gap. No effective useful value of this parameter
is asserted.
\end{theorem}

Although several lemmas below hold for arbitrary fixed $k\geq2$ and
bounded input tuples, the stated transfer target is the source just
specified. No uniformity as $k\to\infty$ is claimed. The parameter $k$
is held fixed in every matrix-size and small-parameter limit.

The contribution is confined to three coupled operations: the
unconditional Fisher bound of Proposition~\ref{prop:fisher}; the
unconditional, all-output-size upper volume bound of
Proposition~\ref{prop:volume}, with the same constant as its lower
bound; and their coupling in Section~\ref{sec:coupling}. The underlying
matrix algebra, orbit volumes, lower packing pattern, orthogonal
semicircular rotation, relative change of variables, and final
entropy-increment comparison are classical or supplied by the source.
Section~\ref{sec:comparison} gives the precise comparison and its
literature-screen limits. In particular, we make no worldwide priority
claim.

\section{Entropy conventions and classical interfaces}
\label{sec:conventions}
Write $\tr_d=d^{-1}\Tr_d$ and equip $M_d(\C)^{\sa}$ with real
inner product $\Tr_d(AB)$. All matrix Lebesgue volumes refer to this
\emph{unnormalized} Hilbert--Schmidt metric. For a $q$-tuple $A$, let
$\Gamma_R(A;m,d,\gamma)$ be the tuples $a\in(M_d^{\sa})^q$ with
$\max_j\|a_j\|\leq R$ and
\[
 |\tr_d(a_{j_1}\cdots a_{j_\ell})-
 \tau(A_{j_1}\cdots A_{j_\ell})|<\gamma
 \quad(1\leq\ell\leq m).
\]
With $\log0=-\infty$, set
\begin{align}
 \chi_R(A)&=\inf_{m\in\N,\,\gamma>0}\limsup_{d\to\infty}
 \left[d^{-2}\log\Vol\Gamma_R(A;m,d,\gamma)+\frac q2\log d\right],
 \label{eq:chi}\\
 \chi(A)&=\sup_{R>0}\chi_R(A).
\end{align}
We use cutoff independence: $\chi_R(A)=\chi(A)$ when
$R>\max_j\|A_j\|$ \cite{Vchi}.

The conjugate system $\xi_j\in L^2(W^*(A),\tau)$ is characterized by
\[
 \tau(\xi_j P(A))=(\tau\otimes\tau)(\partial_jP(A)),\qquad
 \partial_jx_i=\delta_{ij}1\otimes1.
\]
The free difference quotient has its usual outer bimodule Leibniz rule.
Set $\Phi^*(A)=\sum_j\|\xi_j\|_2^2$, or $+\infty$ if no conjugate
system exists. With standard free semicircular noise $S$ free from $A$,
\begin{equation}\label{eq:chistar}
 \chi^*(A)=\frac q2\log(2\pi e)+\frac12\int_0^\infty
 \left(\frac q{1+s}-\Phi^*(A+\sqrt s S)\right)\,ds.
\end{equation}
The standard interfaces used here are the comparison
$\chi\leq\chi^*$ \cite{BCG,JP}, the free heat identity
\begin{equation}\label{eq:heat}
 \chi^*(A+\sqrt t S)-\chi^*(A)
 =\frac12\int_0^t\Phi^*(A+\sqrt s S)\,ds,
\end{equation}
when the quantities are finite, and the projection bound
$\Phi^*(A+\sqrt s S)\leq\Phi^*(A)$ for finite initial Fisher
information \cite{Vstar}. The normalization \eqref{eq:chistar} and the
finite heat identity agree with equations \texttt{eq:chistar} and
\texttt{eq:freeflow} of \cite{source298}. The heat identity follows by
shifting the integral in \eqref{eq:chistar}; the logarithmic reference
terms cancel. Finite initial Fisher information makes its lower
endpoint integrable in our applications.

We also use the covariance upper bound
\begin{equation}\label{eq:covariance}
 \chi^*(A)\leq\frac q2\log(2\pi e c_A),\qquad
 c_A=q^{-1}\sum_j\tau((A_j-\tau(A_j))^2)>0.
\end{equation}
For clarity, the classical derivation is short. Conjugate identities
and Cauchy--Schwarz give
$\Phi^*(A+\sqrt s S)\geq q/(c_A+s)$; if conjugates do not exist the
inequality is automatic. Substitution in \eqref{eq:chistar}, using
$\int_0^\infty((1+s)^{-1}-(c_A+s)^{-1})\,ds=\log c_A$, proves
\eqref{eq:covariance}. A scalar tuple has entropy $-\infty$ and presents
no exceptional finite-entropy case.

The relative entropy $\chi(A\mid B)$ used once below is
Shlyakhtenko's supremum-over-conditioning-models entropy
\cite[Definition 2.3]{Shly}. For a finite conditioning tuple, take the
supremum of the fiber volume over its bounded matrix models, then the
ordinary limsup, infimum over joint moment accuracies, and supremum over
cutoffs; for an algebra take the infimum over finite conditioning
families. It is not defined by subtraction of joint entropies.
The particular classical statements needed are:
\begin{enumerate}[label=(\roman*)]
\item $\chi(A\mid T)\leq\chi(A:T)\leq\chi(A)$, by
\cite[Proposition 2.8]{Shly} and projection of joint microstate sets.
\item If $A$ is free from an embeddable algebra $B$, then
$\chi(A\mid B)=\chi(A)$ \cite[Theorem 2.19]{Shly}.
\item An invertible self-adjoint noncommutative power-series map $F$
with coefficients in $B$, with both inverse series converging on
multiradii strictly exceeding the evaluated norms, satisfies
\begin{equation}\label{eq:relative-cov}
 \chi(F(A)\mid B)=\chi(A\mid B)
 +(\Tr_q\otimes\tau\otimes\tau)(\log|D_BF(A)|).
\end{equation}
This is \cite[Theorem 4.1]{Shly}. The proof is omitted in the read arXiv
v1, by reference to Voiculescu's argument; we use its stated theorem.
Every printed analytic hypothesis is verified in
Section~\ref{sec:straighten}.
\item For embeddable $W^*(A)$ and $t>0$,
\begin{equation}\label{eq:gaussian-lower}
 \chi(A+\sqrt t S)\geq\frac q2\log(2\pi e t).
\end{equation}
This is \cite[Proposition 4.2]{Shly}, with the $q/2$ constant in its
proof. The displayed proposition in arXiv v1 prints $q\log$ instead;
that is a normalization typo. The embeddability hypothesis is essential
to this use.
\end{enumerate}

\section{Uniform normal geometry}
\label{sec:geometry}
For this and the next two sections, fix $k\geq2$ and $n=k^2+1$.
Constants may depend on $k,D$ and the indicated operator-norm bounds,
but never on a matrix approximation size. For tuples, $\|\cdot\|$
denotes the maximum coordinate norm; changing to the sum norm changes
only such constants.

Commuting with $D_1$ forces a scalar matrix to be diagonal, and commuting
with every path edge of $D_2$ equates its diagonal entries. Thus their
common commutant is $\C1$. The real map
\[
 L(H)=(i[H,D_1],i[H,D_2])
\]
has kernel $\R1$ and rank $k^2-1$. Write $\mathcal O=\operatorname{ran}L$
and $\mathcal V=\mathcal O^\perp$, so $\dim_\R\mathcal V=n$.
The basis $V_a$ in \eqref{eq:encoding} is orthonormal for normalized
$\tr_k$. A pair $V$ is normal precisely when
$\sum_s[D_s,V_s]=0$.

At size $kd$ let $\mathcal K_d$ be the orthogonal complement of
$I_k\otimes M_d^{\sa}$ in $M_{kd}^{\sa}$. The map
\begin{equation}\label{eq:linear-chart}
 \mathcal L_d(K,Y)=i[K,D\otimes I_d]+Y,
 \qquad K\in\mathcal K_d,\quad Y\in\mathcal V\otimes_\R M_d^{\sa},
\end{equation}
is an isomorphism onto the pair space. In Hermitian coefficient bases it
is the fixed scalar map tensored with the identity on $M_d^{\sa}$.
Its Hilbert--Schmidt inverse norm is therefore independent of $d$.
Its operator-norm inverse is uniformly bounded as well: extract the
finitely many $k\times k$ coefficients by block maps, apply the fixed
scalar inverse matrix, and reconstruct. Each block map is bounded in
operator norm independently of $d$. This same argument works with
entries in any unital $C^*$-algebra.

\begin{lemma}[Uniform analytic normal tube]\label{lem:tube}
There is $r_0>0$, depending only on $k,D$, such that the associated
bundle map
\begin{equation}\label{eq:bundle}
 [U,Y]\longmapsto U(D\otimes I_d+Y)U^*,\qquad
 U\in U(kd),\quad Y\in\mathcal V\otimes_\R M_d^{\sa},\quad\|Y\|<r_0,
\end{equation}
over $U(kd)/(I_k\otimes U(d))$ is injective, for every $d$.
Locally the normal coordinates and the transverse unitary have
absolutely convergent noncommutative power-series descriptions, with
uniform radii and derivative bounds. These descriptions hold over
any coefficient $C^*$-algebra. On $\|Y\|\leq C_R\eps$, the Jacobian
relative to orbit volume times Euclidean normal-fiber volume lies between
$\exp(-C\eps d^2)$ and $\exp(C\eps d^2)$.
\end{lemma}
\begin{proof}
Let $P_o,P_v$ be the fixed coefficient projections onto the orbit and
normal spaces. Choose the sign of $L$ as above and put
\[
 G(K,W)=P_o\big(e^{-iK}(D+W)e^{iK}-D\big)
       =P_oW-LK+\mathcal R(K,W).
\]
The exponential series gives, for small $K,W$,
\begin{align*}
 \|\mathcal R(K,W)\|&\leq C(\|K\|^2+\|K\|\|W\|),\\
 \|\mathcal R(K,W)-\mathcal R(K',W)\|
 &\leq C(\|K\|+\|K'\|+\|W\|)\|K-K'\|.
\end{align*}
All coefficient projections and $L^{-1}$ are the fixed finite maps
just discussed. Hence
\begin{equation}\label{eq:contraction}
 K=L^{-1}P_oW+L^{-1}\mathcal R(K,W)
\end{equation}
is a strict contraction on
$\|K\|\leq2\|L^{-1}P_o\|\|W\|$ for a sufficiently small uniform
radius.

For the analytic assertion, give a series in finitely many
noncommuting letters the absolute majorant norm obtained by summing
the norms of its coefficients times the products of assigned positive
letter radii. This is a Banach algebra norm and is submultiplicative.
The exponential series and the two estimates above hold in this norm.
Iteration of \eqref{eq:contraction} starting at zero gives an absolutely
convergent series. Choose the radius of actual evaluation strictly
smaller than its majorant radius. Differentiation is then bounded,
since a degree-$\ell$ term acquires at most $\ell$ factors and the
radius margin dominates this factor. The same observation gives a
second derivative bound. The resulting normal displacement is
$P_v(e^{-iK}(D+W)e^{iK}-D)$.

The local map $(K,Y)\mapsto e^{iK}(D+Y)e^{-iK}$ is injective on a
uniform convex small ball. Indeed its derivative differs from
\eqref{eq:linear-chart} by $O(r)$ in operator norm and in
Hilbert--Schmidt operator norm on a ball of radius $r$. The derivative
formula for the exponential and
$\|AB\|_{\HS}\leq\|A\|\|B\|_{\HS}$ prove both bounds.
Integrating $\mathcal L_d^{-1}D\Psi-I$ along a segment gives a
Lipschitz error less than $1/2$ after decreasing $r$.

For global uniqueness, if two normal representatives in
\eqref{eq:bundle} coincide, their transition unitary $u$ satisfies
$\|[u,D_s]\|\leq2r_0$. From $D_1$,
$\|u_{ij}\|\leq2r_0/|i-j|$ for $i\ne j$. In the tree-edge entries
of $[u,D_2]$, the other terms are off-diagonal blocks of this size;
thus $\|u_{ii}-u_{jj}\|=O_k(r_0)$ on each edge and then along every
path. The average diagonal block is close to a unitary, and its polar
part $v$ satisfies $\|u-I_k\otimes v\|=O_k(r_0)$.
Near this subgroup the map $(K,v)\mapsto e^{iK}(I_k\otimes v)$ has a
uniform analytic inverse. To see this without a size-dependent inverse
function theorem, use the bounded coefficient projections onto
$\mathcal K_d$ and its complement, the convergent logarithm near the
identity, and the same contraction argument as \eqref{eq:contraction}.
Removing $v$ reduces to the local injective chart. The representatives
therefore agree in the associated bundle. This argument is valid over
any coefficient algebra, a fact needed when identifying recovered laws.

Finally fix $[U,Y]$ and use coordinates
$(\kappa,y)\mapsto Ue^{i\kappa}(D+Y+y)e^{-i\kappa}U^*$ at that point.
At the coordinate origin the bundle measure has density
$|\det\mathcal L_d|$ relative to the product Lebesgue measure: its
horizontal density is the orbit differential and its vertical density
is Euclidean. The differential of the map is $\mathcal L_d+E$, where
$E(\kappa,y)=i[\kappa,Y]$. Thus its relative Jacobian is
\[
 \frac{|\det(\mathcal L_d+E)|}{|\det\mathcal L_d|}
 =|\det(I+\mathcal L_d^{-1}E)|.
\]
Here $\|\mathcal L_d^{-1}E\|_{\HS\to\HS}=O(\eps)$.
All singular values lie between $1-C\eps$ and $1+C\eps$ in real
dimension $2k^2d^2$, giving the stated exponential bounds. This is a
pointwise local-coordinate calculation, so invariance makes it valid
throughout the bundle.
\end{proof}

\subsection{The common volume constant}
Let $c_D$ be the product of the $k^2-1$ positive singular values of $L$
on the scalar traceless Hermitian space, with normalized trace metrics.
Each is repeated $d^2$ times at size $kd$. The base-orbit volume is
\begin{equation}\label{eq:orbit-volume}
 A_d=c_D^{d^2}\frac{\Vol U(kd)}{k^{d^2/2}\Vol U(d)}.
\end{equation}
The factor $k^{d^2/2}$ comes from the metric induced by
$v\mapsto I_k\otimes v$. The normal coefficient map has Jacobian
\begin{equation}\label{eq:fiber-jac}
 (\eps\sqrt k)^{nd^2}.
\end{equation}
Indeed a Hermitian coefficient basis orthonormal for $\Tr_d$ is
multiplied in pair norm by $\sqrt k$ on tensoring with $V_a$.

The classical unitary volume formula in this metric is
\begin{equation}\label{eq:unitary-volume}
 \Vol U(l)=\frac{(2\pi)^{l(l+1)/2}}{\prod_{j=1}^{l-1}j!},\qquad
 c=\frac34+\frac12\log(2\pi).
\end{equation}
For example, projecting to the first column gives a quotient sphere
whose $2l-2$ horizontal metric directions have squared length multiplied
by $2$, and whose phase direction is unchanged. Its volume is
$(2\pi)^l/(l-1)!$; induction proves the formula. Stirling summation
then yields
\begin{equation}\label{eq:orbit-asymptotic}
 \frac{\log A_d}{d^2}
 =-\frac{k^2-1}{2}\log d-\frac{k^2+1}{2}\log k
 +(k^2-1)c+\log c_D+o(1).
\end{equation}
The orbit/fiber factorization and the constant $c$ are established
antecedents: see Nica--Shlyakhtenko--Speicher
\cite[Lemmas 2.6 and 2.8]{NSSentropy}, including their reference to
Marinov, and the upper/lower orbit patterns in their Lemmas 2.9 and
2.10. We claim no novelty for these constants or the bookkeeping
\eqref{eq:orbit-volume}--\eqref{eq:orbit-asymptotic}.

\section{Matched microstates volumes at every output size}
\label{sec:volume}
\begin{proposition}\label{prop:volume}
Let $A$ be a bounded self-adjoint $n=k^2+1$ tuple with $\chi(A)>-\infty$
and $\max_a\|A_a\|\leq R$. For all sufficiently small $\eps>0$,
\begin{equation}\label{eq:volume}
 \left|\chi(E_\eps(A))-k^{-2}\chi(A)
       -\frac n{k^2}\log\eps-C_D\right|\leq C_{k,R}\eps,
\end{equation}
where
\begin{equation}\label{eq:CD}
 C_D=\log k+\frac{k^2-1}{k^2}c+\frac{\log c_D}{k^2}.
\end{equation}
The constants and smallness threshold can be chosen uniformly for a
fixed coordinate bound $R$. Both sides use \eqref{eq:chi}.
\end{proposition}

\subsection{Lower bound and preservation of source subsequences}
Fix a source cutoff $R_1>\max_a\|A_a\|$. A source microstate set
$\Gamma_{R_1}(A;m,d,\gamma)$ is invariant under simultaneous $U(d)$
conjugation. It therefore defines an associated-bundle fiber.
For any prescribed output moments and a sufficiently large fixed output
cutoff, every point of the bundle over a sufficiently accurate source
microstate set is an output microstate: substitution in
\eqref{eq:encoding} is polynomial, and conjugation preserves every
trace word. Lemma~\ref{lem:tube} and \eqref{eq:fiber-jac} give
Here $\Gamma_{\mathrm{out},kd}$ denotes the prescribed output
microstate neighborhood in size $kd$.
\begin{equation}\label{eq:lower-volume}
 \Vol\Gamma_{\mathrm{out},kd}
 \geq A_d(\eps\sqrt k)^{nd^2}
       \Vol\Gamma_{R_1}(A;m,d,\gamma)e^{-C\eps d^2}.
\end{equation}
Every subsequence realizing a source limsup gives output sizes $kd$.
This suffices for an output lower limsup bound and loses no source
subsequence. Substituting \eqref{eq:orbit-asymptotic} cancels the
$\log d$ terms, since $n+k^2-1=2k^2$. The orbit contribution
$-n\log k/2$ cancels the fiber contribution $n\log k/2$.
After division by $(kd)^2$, the output normalization $\log(kd)$ leaves
\eqref{eq:CD}. Taking infima over moments proves the lower half of
\eqref{eq:volume}. In particular $\chi(E_\eps(A))> -\infty$.

\subsection{Rank perturbations and intrinsic recovery}
The upper bound requires arbitrary output sizes $N$, not just multiples
of $k$. We first give the rank estimate used in recovery.

\begin{lemma}[Rank perturbation of moments]\label{lem:rank}
Suppose two tuples $a,b$ in $M_N$ have coordinate norms at most $R_0\geq1$
and $\rank(a_j-b_j)\leq m$ for every $j$. A word $w$ of length $L$
satisfies
\[
 |\tr_N(w(a)-w(b))|\leq2LR_0^L m/N.
\]
Compression by a projection of corank $m$ changes any fixed trace word
by $O_{L,R_0}(m/N)$, including the normalization change to the retained
corner when $m/N\to0$.
\end{lemma}
\begin{proof}
Telescope $w(a)-w(b)$ into $L$ terms, each containing one difference.
Each term has rank at most $m$ and norm at most $2R_0^L$.
The estimate $|\Tr T|\leq\rank(T)\|T\|$ proves the assertion.
For compression, regard $Pa_jP$ as a matrix on the original space;
$a_j-Pa_jP$ has rank at most $2m$. Apply the first assertion and then
$\tr_{\rank P}=(N/\rank P)\tr_N$ on the corner.
\end{proof}

\begin{lemma}[Uniform recovery from pair microstates]\label{lem:recovery}
Fix $A,R$ as in Proposition~\ref{prop:volume} and a sufficiently small
$\eps>0$. There is $R'>\max_a\|A_a\|$, depending only on $k,R$, with
the following property. Given a finite output cutoff $R_{\mathrm{out}}$,
source parameters $m_0\in\N,\gamma_0>0$, and $\theta>0$, there are
output parameters $m,\gamma$ and $N_0$ such that every
\[
 z\in\Gamma_{R_{\mathrm{out}}}(E_\eps(A);m,N,\gamma),\qquad N\geq N_0,
\]
has a projection $P$ of rank $r=kd\geq(1-\theta)N$ and exact matrix
units on $P\C^N$ for which $PzP$ is in the normal tube and its recovered
coefficients lie in $\Gamma_{R'}(A;m_0,d,\gamma_0)$.
The cutoff $R'$ is independent of $N,\theta,R_{\mathrm{out}}$.
\end{lemma}
\begin{proof}
We prove the sequential assertion first. Let $z^{(N)}$ be any sequence
at a common output cutoff whose joint moments tend to those of
$Z=E_\eps(A)$. All $o(N)$ terms below are along this sequence.

\emph{Spectral bands and their traces.}
The actual $Z_1$ has $k$ spectral bands around $1,\ldots,k$.
Along $D_1+s\eps\sum_aV_{a1}\otimes A_a$, $0\leq s\leq1$, these
bands remain separated for small $\eps$. Their Riesz projections vary
continuously in norm. Projections at distance less than one are
unitarily equivalent, by the polar decomposition of their intertwiner;
subdivision of the path shows that the actual band projections $p_i$
are equivalent to $E_{ii}\otimes1$. Thus
$\tau_{\mathcal N}(p_i)=1/k$ exactly, even when the coefficient algebra
is not a factor.

Choose slightly enlarged band intervals with boundaries in the actual
gaps. The corresponding matrix spectral projections $p_i^{(N)}$ have
ranks $N/k+o(N)$, and their complement has rank $o(N)$. Indeed bounded
moment convergence implies weak convergence of spectral measures;
closed gap neighborhoods have zero limiting mass. Smooth functions
agreeing with the indicators off these neighborhoods differ from the
spectral projections only on rank $o(N)$. Uniform polynomial
approximation then transfers joint moments.

\emph{Tree-polar synchronization.}
In the actual algebra, the projections can be simultaneously identified
with $E_{ii}\otimes1$ by a unitary within $O_{k,R}(\eps)$ of $1$.
For example, take the polar part of $\sum_i p_i(E_{ii}\otimes1)$.
In these coordinates every tree-edge block of $Z_2$ is within
$O_{k,R}(\eps)$ of $1$, so its singular values lie in $[1/2,2]$ after
shrinking $\eps$. Its polar unitary is also close to $1$.

For the matrix edge block $y=p_i^{(N)}z_2p_j^{(N)}$, take its polar
partial isometry on singular values in a fixed, slightly larger
interval, for example $[1/3,3]$. The excluded initial and final ranks
are $o(N)$ by convergence of the spectral distribution of $y^*y$ in
its band corner and the preceding rank approximation. The truncated
polar map is bounded; on the good interval its formula is
$y(y^*y)^{-1/2}$. A continuous extension of this bounded function
approximates it in moments, with discrepancies confined to rank
$o(N)$. No inversion near zero occurs.

Choose band $1$ as root. Along each finite tree path pull the next good
initial domain back to the root and intersect with the domains already
retained. Each step loses at most the excluded rank, together with the
$o(N)$ discrepancy of band ranks. A partial isometry is an isometry
on its good domain, so pulling back does not increase this loss. The
final root subspace has dimension $d_0=N/k-o(N)$. Restricted path
maps $v_i$ are isometries from that root into mutually orthogonal band
subspaces, with $v_1$ the root inclusion. Then
$f_{ij}=v_iv_j^*$ are exact matrix units on their sum.
The absence of cycles removes any consistency condition.

\emph{Root moments.}
The root entries $a_{ij}^{(s,N)}=v_i^*z_sv_j$ are uniformly bounded by
$R_{\mathrm{out}}$. Replace each spectral or polar operation by its
continuous functional-calculus approximation. Every discarded-domain
restriction changes the resulting bounded products by rank $o(N)$.
The root intersection need not itself be a functional calculus:
Lemma~\ref{lem:rank} applies directly to its compression. Consequently
the joint $d_0$-normalized moments of the root entries tend to the
moments of the actual polar-coordinate entries, with the normalized
root trace. Polynomial approximation supplies the continuous-function
step, and $d_0/N\to1/k$ supplies the trace normalization.
In the actual algebra, these entries form a pair $D+O_{k,R}(\eps)$.

\emph{A second root truncation gives norm control.}
For each real and imaginary Hermitian part $h$ of each root entry
$a_{ij}^{(s,N)}-(D_s)_{ij}1$, fix a positive bound
$a_0=C_{k,R}\eps$ for the norm of its actual limiting counterpart.
This bound may be positive even if that counterpart is zero.
A continuous nonnegative function vanishing on $[-a_0,a_0]$ and equal
to one outside $[-2a_0,2a_0]$ shows that the bad spectral subspace of
$h$ has rank $o(N)$. Intersect all the good root subspaces and call
the resulting projection $Q$. There are at most $2k^2$ Hermitian parts
per coordinate, hence at most $4k^2$ in total; this is a fixed number
before $N\to\infty$. Their rank losses therefore sum to $o(N)$.

Restrict \emph{every} path map to this one $Q$. Exact matrix units
remain, now on a corner of rank $kd=N-o(N)$. Since
$Q\leq1_{[-2a_0,2a_0]}(h)$,
$\|QhQ\|\leq2a_0$; invariance of $Q$ under $h$ is unnecessary.
Each complex entry deviation has norm at most $4a_0$, so a conservative
bound for the whole pair is $4k^2a_0=O_{k,R}(\eps)$.
Recompression changes each prescribed moment by $o(1)$ by
Lemma~\ref{lem:rank}.

\emph{Normal inversion and source-law identification.}
Decrease $\eps$ so that this last norm bound is strictly inside the
analytic domain of Lemma~\ref{lem:tube}. The normal inverse gives
coefficient norms at most $C_k a_0/\eps\leq C'_{k,R}$.
Choose $R'$ strictly larger than that bound and $\max\|A_a\|$.
The actual polar-coordinate pair is a unitary conjugate of $E_\eps(A)$
by a unitary close to the scalar stabilizer. The uniqueness in
Lemma~\ref{lem:tube}, valid over the actual coefficient algebra,
identifies its recovered coefficients with a common unitary conjugate
of $A$. Hence their joint law is that of $A$.
The inverse series converges absolutely on a ball strictly containing
both actual and truncated matrix entries. Uniform polynomial
approximation to this series transfers every desired source moment.
This proves the sequential recovery assertion and its fixed cutoff.

Finally, if the quantified statement failed for fixed
$R_{\mathrm{out}},m_0,\gamma_0,\theta$, one could choose bad microstates
with output accuracies $m_j\to\infty$, $\gamma_j\to0$ and sizes
$N_j\to\infty$. They form a bounded moment-convergent sequence,
contradicting the sequential assertion. This proves uniformity over
\emph{all} sufficiently accurate output microstates.
\end{proof}

\subsection{The discarded-block volume and the ordinary limsup}
We complete the upper bound of Proposition~\ref{prop:volume}.
Fix $\eps$, an output cutoff, and a desired source neighborhood, leaving
strict slack in its moment tolerances and cutoff. Let $r=kd$, $l=N-r$.
Rank-$l$ projections admit an operator-norm $\eta$-net of cardinality
\begin{equation}\label{eq:net}
 \#\mathcal P_{N,l,\eta}\leq(C/\eta)^{2Nl}.
\end{equation}
To prove this, represent them by $N$-by-$l$ isometric frames, all in
the operator-norm unit ball of a real $2Nl$-dimensional normed space.
A maximal separated set has the usual disjoint-ball volume bound;
passing to projections changes only $C$.

For projections with $\|P-Q\|<\eta<1/4$, the polar part $w$ of
$QP+(1-Q)(1-P)$ satisfies $wPw^*=Q$ and $\|w-1\|\leq C\eta$.
Thus $QzQ$ differs in norm by at most $CR_{\mathrm{out}}\eta$ from
the transported recovered block $wPzPw^*$. The analytic inverse is
locally Lipschitz, with constants allowed to depend on fixed $\eps$;
choose $\eta$ small enough that the recovered source coefficients of
$QzQ$ lie in the prescribed relaxed source neighborhood. Strict slack
keeps it inside the normal tube. Because the tube image is
$U(r)$-invariant, no choice of frame for $Q$ matters. This covers the
union of possible recovered corners without choosing a measurable
recovery map and without estimating an uncontrolled thickening of a
source microstate set.

For each net corner, Lemma~\ref{lem:tube} and the area formula bound
its pair-block volume by
\begin{equation}\label{eq:upper-corner}
 A_d(\eps\sqrt k)^{nd^2}\Vol\Gamma_{R'}(A;m_0,d,\gamma_0)
 e^{C\eps d^2},
\end{equation}
where the neighborhood on the right is the relaxed one.
The remaining Hermitian pair coordinates have real dimension
$M=2(N^2-r^2)$ and lie in a Euclidean ball of radius
$R_{\mathrm{out}}\sqrt{2N}$. Orthogonal decomposition of the ambient
Hilbert--Schmidt space makes Lebesgue measure factor into the corner
and these remaining coordinates. Cross blocks carry their induced
factor-two metric; that metric is already included in this ball.
The Euclidean ball formula and Stirling's bound give, when $M>0$,
\begin{align}
 \log\Vol_{\rm rem}
 &\leq-(N^2-r^2)\log N\notag\\
 &\quad+C(N^2-r^2)
 \left[1+\log^+R_{\mathrm{out}}+
                  \log\frac{N^2}{N^2-r^2}\right].
 \label{eq:remainder-ball}
\end{align}
For $M=0$ the complementary volume is $1$.

Multiply \eqref{eq:upper-corner} by this bound and \eqref{eq:net},
then sum over the at most $N+1$ possible ranks. Let
$\theta'=1-r/N\leq\theta$. The source logarithm, orbit logarithm,
and fiber factor combine into
\[
 -r^2\log r+r^2
 \left[k^{-2}h_d+\frac n{k^2}\log\eps+C_D+O(\eps)+o(1)\right],
\]
where
$h_d=d^{-2}\log\Vol\Gamma_{R'}(A;m_0,d,\gamma_0)+(n/2)\log d$.
The logarithmic size terms, including the output normalization, are
\[
 -r^2\log r-(N^2-r^2)\log N+N^2\log N
 =r^2\log(N/r).
\]
After division by $N^2$, all complementary, net, rank-sum, and
normalization errors are bounded by
\begin{equation}\label{eq:theta-error}
 C\left[\theta\log(C/\theta)+\theta\log(C/\eta)\right]+o(1).
\end{equation}
Constants here may depend on the fixed output cutoff and $\eps$.

For fixed source parameters, the definition of a limsup bounds $h_d$
above for every sufficiently large $d$ by its limsup plus an arbitrary
positive error. If that limsup is $-\infty$, use an arbitrary finite
upper bound instead and then send it to $-\infty$. Every retained
$d\geq(1-\theta)N/k$ tends to infinity; thus this estimate applies to
all output sizes, regardless of their residue modulo $k$.
Lemma~\ref{lem:recovery} allows $\theta$ to be arbitrarily small by
choosing the output moment accuracy sufficiently strong.

Explicitly, hold $\eps$, source accuracy, and output cutoff fixed;
choose $\eta$ for the source slack; choose output accuracy for a given
$\theta$; take the limsup as $N\to\infty$; let $\theta\downarrow0$;
then let the net slack and source tolerance tend to zero and the source
moment order to infinity. Lastly take the supremum over output
cutoffs. The surviving constant $C_D$ and $O_{k,R}(\eps)$ error are
independent of the output cutoff, since the recovered cutoff $R'$ is.
Cutoff independence gives $\chi_{R'}(A)=\chi(A)$. This proves the
upper half of \eqref{eq:volume}, and completes its proof.

\section{Unconditional scores in the coarse bimodule}
\label{sec:fisher}
This section supplies the estimate that cannot be obtained by simply
conditioning on scalar matrix units. Work in
$\mathcal N=M_k(\mathcal A)$ with normalized trace $\tr_k\otimes\tau$.
The coarse bimodule has outer actions
$x\cdot(a\otimes b)\cdot y=xa\otimes by$. Inner products use normalized
traces. Tensor notation may equivalently use the opposite algebra on
the second factor; the displayed outer-action convention fixes all
matrix-unit orientations.

Let $O_1,\ldots,O_{k^2-1}$ be a real orthonormal basis of $\mathcal O$.
For bounded self-adjoint tuples $U=(U_a)_{a=1}^n$ and
$T=(T_b)_{b=1}^{k^2-1}$, put
\begin{equation}\label{eq:noisy-form}
 Z_s=D_s+\eps\left(\sum_aV_{as}\otimes U_a+
                         \sum_bO_{bs}\otimes T_b\right).
\end{equation}

\begin{proposition}\label{prop:fisher}
Suppose $U$ has finite conjugate system $\xi_a$ and is free from $T$.
For a common coordinate norm bound $R$, and sufficiently small
$\eps>0$, the pair \eqref{eq:noisy-form} has unconditional conjugates
and
\begin{equation}\label{eq:fisher}
 \Phi^*(Z)\leq\frac{\Phi^*(U)}{\eps^2k^4}
       +C_{k,R}\eps^{-1}(1+\Phi^*(U)).
\end{equation}
\end{proposition}

\subsection{The finite coarse decomposition}
In $M_k\otimes M_k$ set
\[
 w_{bc}=\sum_{i=1}^k E_{ib}\otimes E_{ci},\qquad
 W_{a,bc}=(V_{a1}w_{bc},V_{a2}w_{bc}).
\]
For every scalar matrix $H$, $Hw_{bc}=w_{bc}H$, by direct
matrix-unit multiplication. Conversely, solving these equations for
all $E_{ij}$ shows that the $w_{bc}$ span the scalar-central subspace.
The same calculation, with normalized tensor traces, gives
\begin{equation}\label{eq:gram}
 \langle W_{a,bc},W_{e,fg}\rangle
 =k^{-1}\delta_{ae}\delta_{bf}\delta_{cg}.
\end{equation}
Let $J\eta=([D_1,\eta],[D_2,\eta])$ in the coarse pair space.
Since $\sum_s[D_s,V_{as}]=0$, $J^*W_{a,bc}=0$.
The kernel of $J$ is the scalar-central subspace, of complex dimension
$k^2$, because $D_1,D_2$ generate $M_k$. Hence
\[
 \rank J=k^4-k^2,\qquad
 \dim\operatorname{span}\{W_{a,bc}\}=k^2(k^2+1).
\]
These dimensions add to $2k^4$. Thus $J$ on its kernel complement and
the span of the $W$ give an orthogonal direct-sum isomorphism of the
entire coarse pair space. This is a calculation in the \emph{coarse}
module, rather than an inference from the ordinary tangent count.
Tensoring the finite coefficient identity with
$\mathcal A\otimes\mathcal A$ preserves it.

\subsection{Derivations, coefficient placement, and bilinear scores}
Because $U$ is free from $T$, its conjugates $\xi_a$ are also conjugates
relative to $\alg(T)$. Existence of these conjugates makes the free
difference quotients well defined on the represented algebra: the
adjoint formula against polynomial elementary tensors makes the
derivative of a zero polynomial orthogonal to a dense set, hence zero.
This standard domain implication is part of the conjugate-variable
calculus \cite{Vstar}.

Define entry derivations $\Delta_{a,bc}$, vanishing on $M_k$ and on
$T$, with
$\Delta_{a,bc}(U_e)=\delta_{ae}w_{bc}$.
Compatibility with scalar-matrix relations follows from centrality of
$w_{bc}$. On a matrix monomial $E_{ij}\otimes P(U,T)$, trace
contraction of this derivative is
$\delta_{ib}\delta_{jc}/k^2$ times the entry difference-quotient
contraction. Throughout this section a \emph{bilinear score}
$\mathfrak b(\delta)$ means a vector satisfying
\[
 (\tr_k\otimes\tau)(\mathfrak b(\delta)P)
 =((\tr_k\otimes\tau)\otimes(\tr_k\otimes\tau))(\delta(P))
\]
for every represented polynomial $P$. This is not Hilbert-adjoint
notation. For the Hilbert inner product $\langle x,y\rangle=\tau(x^*y)$,
a Hilbert adjoint at $1\otimes1$ gives
$\mathfrak b(\delta)=(\delta_H^*(1\otimes1))^*$.
Trace contraction of the entry derivation gives
\begin{equation}\label{eq:entry-adjoint}
 \mathfrak b(\Delta_{a,bc})=k^{-1}E_{cb}\otimes\xi_a.
\end{equation}
The transpose and $k^{-1}$ in this formula are essential.
It is the same entry-adjoint method underlying
\cite[Proposition 4.1]{NSSfisher}; the normalization here uses real
Hermitian coefficients orthonormal for $\tr_k$.

For a coefficient $u\otimes v\in\mathcal A\otimes\mathcal A$,
insert it on the \emph{inner} legs of the derivative: the contribution
at an occurrence $P U_a Q$ is
\begin{equation}\label{eq:insertion}
 \sum_i P(E_{ib}\otimes u)\ \otimes\ (E_{ci}\otimes v)Q.
\end{equation}
This insertion commutes with both outer actions and therefore preserves
the Leibniz rule. Arbitrary outer multiplication of an already formed
derivative would not do so and is not used. Include also inner
derivations $P\mapsto[P,\eta]$.

For precision, for polynomial coefficients $u,v\in\mathcal A$ define
\begin{equation}\label{eq:adjoint-product}
 \mathcal C_a(u,v)=v\xi_a u
 -v(\tau\otimes\id)(\partial_a u)
 -(\id\otimes\tau)(\partial_a v)u.
\end{equation}
If $\delta_{uv}(P)=\sum_{P=P_1U_aP_2}P_1u\otimes vP_2$, then
\begin{equation}\label{eq:bilinear-pairing}
 (\tau\otimes\tau)(\delta_{uv}(P))
 =\tau(\mathcal C_a(u,v)P).
\end{equation}
Indeed apply the conjugate relation to $uPv$, expand
\[
 \partial_a(uPv)=(\partial_a u)(1\otimes Pv)
 +(u\otimes1)(\partial_aP)(1\otimes v)
 +(uP\otimes1)(\partial_a v),
\]
and move factors cyclically inside each scalar trace. The middle term
has trace $\sum\tau(P_1u)\tau(vP_2)$. The other two terms give the two
subtractions in \eqref{eq:adjoint-product}, while
$\tau(\xi_a uPv)=\tau(v\xi_a uP)$. This proves the identity without
mixing bilinear and sesquilinear conventions. For the matrix insertion
\eqref{eq:insertion}, the bilinear score is
$k^{-1}E_{cb}\otimes\mathcal C_a(u,v)$, by the same entry contraction
as in \eqref{eq:entry-adjoint}.

The resulting bound is
\[
 \|\mathcal C_a(u,v)\|_2\leq
 \|u\|\|v\|\|\xi_a\|_2
 +\|v\|\|\partial_a u\|_\pi
 +\|\partial_a v\|_\pi\|u\|.
\]
Thus a sum of inserted coefficients has the same projective coefficient
and differentiated-coefficient bounds used below. The reversed order
of factors changes no norm bound. For an inner derivation, trace
contraction gives the difference of the two partial traces of $\eta$;
its bilinear score has norm at most $2\|\eta\|_\pi$. Equivalently these
are Hilbert-adjoint bounds after the explicit adjoint conversion above.

\subsection{Exact inverse and the leading energy}
Let $H=VU+OT$ denote the bounded perturbation in \eqref{eq:noisy-form}.
For a sum of coefficient-inserted normal derivatives and an inner
derivation, prescribe its two values on $Z$. Use unknowns
$(\eps c,\eta)$, where $c$ collects all normal coefficients and
$\eta$ lies in the fixed complement of $\ker J$. The map to generator
values is
\begin{equation}\label{eq:score-inverse}
 \mathcal L_0+\eps\mathcal B_H.
\end{equation}
Here $\mathcal L_0$ is exactly the finite direct-sum isomorphism from
the preceding subsection. The normal derivations vanish on $T$ and
$D$, and their values on $VU$ are constant, so the only perturbation is
$\eps[H,\eta]$.

Equip the finitely many coefficient arrays with sum projective tensor
norm. If $K=\|\mathcal L_0^{-1}\|$ and
$\|\mathcal B_H\|\leq H_R$, choose $\eps KH_R\leq1/2$.
Then the Neumann series for \eqref{eq:score-inverse} converges absolutely,
has norm at most $2K$, and differs from $\mathcal L_0^{-1}$ by at most
$2\eps K^2H_R$. Its entries are series of words in $U,T$ on the two
legs. Differentiating a degree-$\ell$ summand with respect to an entry
$U_a$ inserts a derivative at at most $\ell$ positions. The
nonconstant part of the inverse therefore has differentiated norm
bounded by a constant times
\[
 \sum_{\ell\geq1}\ell\eps(\eps KH_R)^{\ell-1}=O_{k,R}(\eps).
\]
After undoing the scaling $\eps c$, the correction to the leading
normal coefficient and its differentiated norm are $O_{k,R}(1)$;
the inner coefficient and its norm are also $O_{k,R}(1)$.
Apply \eqref{eq:adjoint-product}--\eqref{eq:bilinear-pairing} first
to every finite Neumann sum, with the inner-derivation score just given.
The coefficient and differentiated-coefficient series converge
absolutely, so their bilinear scores are Cauchy in $L^2$. For a fixed
polynomial $P$, the inserted derivatives converge as well; passing to
the limit in the trace pairing gives a bilinear score for the limiting
derivation. The corresponding Hilbert adjoint exists by the stated
adjoint conversion. This supplies the same remainder norm bound.

Apply the inverse to the target pair
$(\delta_{js}1\otimes1)_{s=1}^2$ to obtain a derivation $\delta_j$
with $\delta_j(Z_s)=\delta_{js}1\otimes1$. It agrees on every
polynomial in $Z$ with its free difference quotient, by Leibniz.
Project its bilinear score onto $L^2(W^*(Z))$. The trace pairing
against every polynomial in $Z$ is preserved by this conditional
expectation, so the projection is the unconditional conjugate variable
of $Z_j$.

The projection of the target pair onto the normal coarse subspace,
computed from \eqref{eq:gram}, is
\[
 \frac1k\sum_{a,b,c}(V_{aj})_{cb}W_{a,bc}.
\]
Using \eqref{eq:entry-adjoint} gives the leading unprojected score
\begin{equation}\label{eq:leading-score}
 \Xi_j^0=\frac1{\eps k^2}\sum_a V_{aj}\otimes\xi_a,
 \qquad
 \sum_{j=1}^2\|\Xi_j^0\|_2^2
 =\frac{\Phi^*(U)}{\eps^2k^4}.
\end{equation}
Indeed the cross terms sum to
$\sum_j\tr_k(V_{aj}V_{ej})=\delta_{ae}$; no orthogonality among the
$\xi_a$ is assumed. The remainder estimate just proved is
\[
 \left(\sum_j\|\Xi_j-\Xi_j^0\|_2^2\right)^{1/2}
 \leq C_{k,R}(1+\sqrt{\Phi^*(U)}).
\]
Squaring the triangle inequality and using contraction of orthogonal
projection proves \eqref{eq:fisher}, absorbing the bounded square
remainder into $C\eps^{-1}(1+\Phi^*(U))$ for $\eps\leq1$.
This completes the proof of Proposition~\ref{prop:fisher}.
The inner derivations supply the orbit component of the target;
the normal leading scores alone would fail the generator tests.

\section{The exact noisy path and analytic straightening}
\label{sec:coupling}
Return to the source $X$ of \eqref{eq:source}, and write
$Z_\eps=E_\eps(X)$. Let $S^{\rm pair}$ be a standard free
semicircular pair free from the full ambient algebra.

\subsection{Noise normalization and the nonmicrostates increment}
A standard scalar semicircular pair over $M_k(\mathcal A)$ can be
realized by $2k^2$ standard free semicircular entry coefficients divided
by $k$, in an orthonormal Hermitian pair basis for normalized $\tr_k$.
This is the classical free semicircular/circular matrix construction
\cite{Vmatrix}. One way to check the normalization is the completeness
identity $\sum_\alpha B_\alpha A B_\alpha=k\Tr_k(A)I$ for a
$\tr_k$-orthonormal Hermitian basis. After division by $k^2$ the
operator-valued covariance is the scalar normalized trace; higher
free cumulants vanish. The two pair-coordinate covariances are
orthogonal, yielding a standard free pair free from $M_k(\mathcal A)$.

Orthogonal rotation to the bases $V$ and $O$ preserves a standard free
semicircular family. Thus, exactly in distribution at
$t=\eps^2k^2u$,
\begin{equation}\label{eq:noise-split}
 Z_\eps+\sqrt t S^{\rm pair}
 =D+\eps\big(VU(u)+OT(u)\big),\quad
 U(u)=X+\sqrt u S^{\rm n},\quad T(u)=\sqrt u S^{\rm o}.
\end{equation}
All coefficient semicirculars are free from $X$, and the normal and
orbital families are free from one another. In particular $U(u)$ is
free from $T(u)$. This is an exact linear identity of laws, so there is
no use of upper semicontinuity of Fisher information or of a separate
hidden copy of $X$.

Fix $U_0<\infty$. The coefficient norms in \eqref{eq:noise-split} are
uniformly bounded for $0\leq u\leq U_0$, and
$\Phi^*(U(u))\leq\Phi^*(X)<\infty$. Proposition~\ref{prop:volume}
and $\chi\leq\chi^*$, together with \eqref{eq:covariance}, show that
both entropies of $Z_\eps$ are finite. For its noisy endpoints the
Fisher bound is integrable on every such interval, and the standard
heat identity applies. Proposition~\ref{prop:fisher} and
$dt=\eps^2k^2\,du$ give
\begin{align}
 &\chi^*(Z_\eps+\eps k\sqrt{U_0}S^{\rm pair})-\chi^*(Z_\eps)
 \notag\\
 &\quad\leq\frac1{2k^2}\int_0^{U_0}\Phi^*(U(u))\,du
                +C_{U_0}\eps(1+\Phi^*(X))\notag\\
 &\quad=k^{-2}\big(\chi^*(X+\sqrt{U_0}S^{\rm n})-\chi^*(X)\big)
                +O_{U_0}(\eps).
 \label{eq:star-increment}
\end{align}
In particular the leading $k^{-4}$ energy in
\eqref{eq:leading-score} becomes $k^{-2}$ after changing heat time.
The error constant can depend on $U_0$, which is fixed before
$\eps\downarrow0$.

\subsection{The precise power-series hypotheses}
\label{sec:straighten}
Fix a positive $U_0$ and abbreviate $U=U(U_0)$, $T=T(U_0)$.
Apply the normal solver \eqref{eq:contraction} to
$W=\eps(VU+OT)$. Its first-order term $K_1=L^{-1}P_oW$ depends only
on $T$, since $P_oVU=0$. Consequently
\begin{equation}\label{eq:straighten}
 e^{-iK}(D+\eps(VU+OT))e^{iK}
 =D+\eps V H_\eps(U,T),\qquad
 H_\eps(U,T)=U+\eps R_\eps(U,T).
\end{equation}
The remainder and its derivative in $U$ have bounded absolute
power-series majorant norms on fixed radii strictly larger than the
coefficient norms. The cancellation of the first-order orbital part
leaves the normal part $\eps VU+O(\eps^2)$, so division by $\eps$
causes no singular remainder.

Here are the hypotheses of \cite[Theorem 4.1]{Shly} in its own terms.
Let $B=\alg(T)$ be the unital algebra generated by the orbital tuple;
its von Neumann closure is embeddable. Group consecutive $T$-letters
in each monomial of the series in $(U,T)$. This expresses
$F=H_\eps$ as a noncommutative power series in the $U$ letters with
coefficients in $B$. The sum of products of coefficient norms is
bounded by the original majorant evaluated at a $T$-radius exceeding
$\|T\|$. Thus it has a multiradius strictly exceeding every $\|U_a\|$.

The inverse series is constructed, with a radius margin, by the
contraction
\begin{equation}\label{eq:inverse-H}
 A=H-\eps R_\eps(A,T).
\end{equation}
Choose a coefficient bound larger than both the evaluated $U$ and $H$
norms, then a strictly larger majorant radius. Uniform bounds on the
remainder and its derivative make \eqref{eq:inverse-H} a contraction
in the associated Banach algebra of convergent series for small
$\eps$. Iteration yields $G$; uniqueness proves both formal identities
$F\circ G=\id$ and $G\circ F=\id$, as well as their analytic
identities. The radii for $F$ and $G$ can be chosen with a common
strict margin beyond the evaluated norms, as required by the printed
condition $\|U_a\|<\min(r_a,R_a)$.

Both $F$ and $G$ preserve self-adjointness. The normal solver operates
in the self-adjoint part, $V$ is a real Hermitian basis, and coefficient
extraction is the scalar partial trace
$\sum_s(\tr_k\otimes\id)(V_{as}Y_s)$. This is self-adjoint for
Hermitian inputs because the scalar matrix factor can be moved through
the partial trace. The same holds for the fixed point inverse.

Finally the coefficient-algebra Jacobian satisfies
\[
 D_BF(U)=I+O_{U_0}(\eps)
 \quad\hbox{in }M_n(\mathcal A\,\overline\otimes\,\mathcal A^{\op}).
\]
The projective majorant estimate implies this operator-norm estimate.
The singular values lie between $1-C\eps$ and $1+C\eps$, hence
\begin{equation}\label{eq:logjac}
 \big|(\Tr_n\otimes\tau\otimes\tau)\log|D_BF(U)|\big|
 \leq C_{U_0}n\eps.
\end{equation}
Since $n$ is fixed it is absorbed into $C_{U_0}$.
This verifies the convergent inverse, strict multiradius,
self-adjointness, and logarithmic Jacobian requirements; a merely
formal inverse would not suffice.

\subsection{The microstates increment}
Theorem 2.19 of \cite{Shly} applies because $U$ is free from $B$ and
$W^*(B)$ is the embeddable semicircular algebra. In this particular
application $T$ has bounded matrix models at every sufficiently large
size. Uniform asymptotic freeness of bounded deterministic $U$ models
with independent GUE models for $T$, followed by averaging over $T$,
provides a conditioning model whose good $U$-fiber has proportion
tending to one. This explains directly why free conditioning retains
every source limsup subsequence here; it does not require aligning two
unrelated maximizing subsequences.

Using Theorems 2.19 and 4.1, then Proposition 2.8 and projection of
joint microstate sets, yields
\begin{equation}\label{eq:H-entropy}
 \chi(H_\eps(U,T))\geq\chi(H_\eps(U,T)\mid B)
 =\chi(U)+(\Tr_n\otimes\tau\otimes\tau)\log|D_BF(U)|
 \geq\chi(U)-C_{U_0}\eps.
\end{equation}
This inequality does not use the equality left to the reader in
Proposition 2.12 of \cite{Shly}. Also $\chi(U)>-\infty$ by
\eqref{eq:gaussian-lower}, since $X$ is embeddable.
The coefficient norm bound for $H_\eps$ is uniform for small $\eps$.

Unitary conjugation in the tracial algebra preserves the joint law,
so \eqref{eq:straighten} and the lower half of
Proposition~\ref{prop:volume} give a lower bound at the noisy endpoint.
Use its upper half at $X$. Their identical constants $C_D$ and
$(n/k^2)\log\eps$ cancel, giving
\begin{align}
 &\chi(Z_\eps+\eps k\sqrt{U_0}S^{\rm pair})-\chi(Z_\eps)
 \notag\\
 &\qquad\geq k^{-2}\big(\chi(X+\sqrt{U_0}S^{\rm n})-\chi(X)\big)
                   -O_{U_0}(\eps).
 \label{eq:micro-increment}
\end{align}
Only a relative \emph{microstates} change-of-variables theorem was
used. The nonmicrostates estimate is the separate unconditional
construction in Section~\ref{sec:fisher}.

\section{Gap transfer and finiteness}
\label{sec:finish}
Let $g=\gap(X)\geq1/2$ and $X_u=X+\sqrt u S^{\rm n}$.
Every entropy in \eqref{eq:star-increment} and
\eqref{eq:micro-increment} is finite: the lower bounds just proved
supply finite microstates entropy, while comparison and covariance
bound the nonmicrostates entropy below and above.
For the noisy pair, $\chi\leq\chi^*$; therefore subtraction of the
two finite increment inequalities gives
\begin{equation}\label{eq:gap-bound}
 \gap(Z_\eps)\geq k^{-2}\big(g-\gap(X_{U_0})\big)-C_{U_0}\eps.
\end{equation}
This final comparison is the standard two-noise-scale strictness
pattern used in \cite{source298}. The work here supplies the two
unconditional estimates to which that pattern is applied.

Embeddability and \eqref{eq:gaussian-lower} give
$\chi(X_u)\geq(n/2)\log(2\pi e u)$ for $u>0$. If
$C_X=n^{-1}\sum_a\operatorname{Var}(X_a)$, covariance and comparison
give
\begin{equation}\label{eq:large-noise-gap}
 0\leq\gap(X_u)\leq\frac n2\log\frac{u+C_X}{u}\longrightarrow0.
\end{equation}
For each fixed $U_0$ let $\eps\downarrow0$ in
\eqref{eq:gap-bound}, and only then let $U_0\to\infty$.
This proves \eqref{eq:transfer}.
For an explicit order of existential choices, first choose $U_0$ so
that $\gap(X_{U_0})<g/2$; then choose $\eps>0$ below all geometric
and analytic thresholds and so small that
$C_{U_0}\eps<g/(4k^2)$. Equation \eqref{eq:gap-bound} then gives
$\gap(Z_\eps)>g/(4k^2)>0$.
Proposition~\ref{prop:volume} gives $\chi(Z_\eps)>-\infty$, and
\eqref{eq:covariance} gives $\chi^*(Z_\eps)<\infty$.
The tensor trace on $M_k(W^*(X))$ is faithful and normal, and both
coordinates of \eqref{eq:encoding} are bounded and self-adjoint.
This proves Theorems~\ref{thm:transfer} and \ref{thm:main}.

Every matrix-size limit in Section~\ref{sec:volume} precedes this
$\eps$ limit. In particular no uniformity of microstate neighborhoods
as $\eps\downarrow0$, no replacement of a limsup by a liminf in
\eqref{eq:chi}, and no ultrafilter are part of the argument.

\section{Relation to prior operations and scope of comparison}
\label{sec:comparison}
\subsection{Orbit volumes and supplied matrix-unit compression}
The polar measure factorization of \cite[Lemma 2.6]{NSSentropy},
its unitary-volume asymptotic in Lemma 2.8, and its upper/lower volume
patterns in Lemmas 2.9 and 2.10 are antecedents of the geometric
bookkeeping here. The last lower-packing lemma has free Haar and
positive-variable hypotheses. Neither that special-case theorem nor
the classical formula \eqref{eq:unitary-volume} by itself proves the
intrinsic recovery statement of Lemma~\ref{lem:recovery}. Uniform
normal-tube geometry is supporting analysis, not a separate novelty
claim. The specific additional upper-volume operation is recovery
from the pair's own microstates with a common-root truncation and
vanishing-rank accounting, valid for every $N$ with the ordinary
limsup and the same $C_D$.

Shlyakhtenko \cite[Section 5, Theorem 5.1]{Shly} already uses the
corner/compression/packing pattern. Its proof takes a corner
$p\leq e_{11}$ with $[e_{11},p]=0$ under \emph{supplied} matrix-unit
conditioning models, compresses with the map $T$ made from those
models, and packs entry microstates back by a measure-preserving map.
This pattern is credited here. Its equality is an ultrafilter entropy
identity, with its specified dilated ultrafilter. The displayed
ordinary-limsup statement is one-sided, and its proof is explicitly
left to the reader. The second displayed intermediate claim repeats
the first inequality direction; the following packing proof establishes
the reverse direction, a printed direction error. We use neither that
misprinted direction nor an inferred ordinary-limsup equality.
The upper bound here instead recovers its matrix units intrinsically
and handles all output ranks.

Nica--Shlyakhtenko--Speicher \cite[Proposition 4.1]{NSSfisher} proves
matrix-entry Fisher identities relative to the scalar matrix algebra,
and unconditional projection inequalities. Its raw complex-entry
normalization has a factor $k^3$. Formula \eqref{eq:leading-score}
uses real Hermitian coefficients normalized for $\tr_k$, yielding
$k^4$ in the score energy. The classical entry-adjoint method and
orthogonal Gaussian rotation are used openly. The additional
unconditional estimate combines normal entry derivatives with inner
derivatives and controls their adjoints; it is not obtained by deleting
matrix-unit conditioning from a relative identity.

\subsection{Why relative identities do not prove the finite gap}
At $Z_\eps=D+\eps VX$, its $2k^2$ real entry coordinates lie in an
affine subspace of dimension $k^2+1$. Thus there are $k^2-1$
independent real-linear relations. Entry microstates entropy is
$-\infty$, including each ultrafilter version. Under the ultrafilter
identity in \cite[Theorem 5.1]{Shly}, the corresponding relative
matrix microstates entropy is $-\infty$. The one-sided
ordinary-limsup statement, whose lower side is then $-\infty$, does
\emph{not} compute the ordinary-limsup relative value.

On the Fisher side, the paid relative matrix-entry identity on the
positive noisy path has normal and orbital contribution
\[
 \Phi^*(Z_\eps+\sqrt t S^{\rm pair}:M_k)
 =\frac{\Phi^*(U(u))+(k^2-1)/u}{\eps^2k^4},
 \qquad t=\eps^2k^2u.
\]
Its orbital term is $((k^2-1)/k^2)t^{-1}$, nonintegrable at zero.
The relative nonmicrostates entropy at the unperturbed pair is
therefore $-\infty$. Between positive source noise times $u_0$ and
$U_0$, this orbital term contributes
\[
 \frac{k^2-1}{2k^2}\log(U_0/u_0).
\]
Thus the relative identities at the target are degenerate. Where the
relative quantities are $-\infty$ and the unconditional ones finite,
conditioning-removal excesses are $+\infty$. Subtracting such excesses
is not a proof of a finite unconditional gap. In the ordinary-limsup
microstates case the cited one-sided theorem leaves even that relative
value undetermined. Our proof never subtracts infinities; it compares
direct unconditional finite increments.

\subsection{Generation, source scope, and limits of the screen}
The pinned preprint \cite{generation}, Theorem \texttt{thm:main},
gives two self-adjoint $W^*$-generators for each separable-predual
$\mathrm{II}_1$ factor. Its introduction credits the classical
$2\times2$ matrix-generation techniques of Douglas--Pearcy and Wogen,
using Dykema--Sinclair--Smith--White Proposition 1.1. Its own
Corollary \texttt{cor:entropy-generators} shows that the free entropy
dimensions $\delta,\delta_0$ depend on the generating tuple.
That theorem supplies generation only; it supplies no finite
$\chi/\chi^*$-gap invariance under replacement of generators.
The companion \cite{source287}, in its amplification and
entropy-dimension sections, transports tuples coordinatewise through
trace-preserving isomorphisms. This preserves their joint law and
length. It also does not reduce the number of entropy coordinates.
We assert no generation property of $Z_\eps$ and use no
factor-isomorphism invariance of entropy.

The source \cite{source298} supplies the many-variable finite gap,
finite initial Fisher information, and bounded matrix models. The
Gaussian lower bound, covariance maximum, comparison $\chi\leq\chi^*$,
relative analytic change of variables, and two-scale strictness
pattern are all credited inputs or classical deductions. The claimed
additional mechanism is precisely the unconditional Fisher bound,
the matched all-size upper volume bound, and their coupling.

The comparison is bounded. Operative bodies of the three primary
papers \cite{Shly,NSSentropy,NSSfisher} and the specified pinned source
and companion statements were read. Dimension, orbital-free-entropy,
and free-mutual-information literature remains a further screen,
including Jung's \href{https://arxiv.org/abs/math/0410594}{math/0410594},
Shlyakhtenko's \href{https://arxiv.org/abs/2201.09503}{2201.09503},
Hiai--Miyamoto--Ueda's
\href{https://arxiv.org/abs/math/0702745}{math/0702745}, Ueda's
\href{https://arxiv.org/abs/1210.6421}{1210.6421} and
\href{https://arxiv.org/abs/1610.04118}{1610.04118},
Izumi--Ueda's \href{https://arxiv.org/abs/1306.5372}{1306.5372}, and the
Connes--Shlyakhtenko framework of dimensions and inner derivations. Their
operative theorem bodies were not screened for this manuscript, so
no claim that these papers lack a relevant transfer theorem is made.
This limitation concerns literature priority; none is invoked as a
mathematical premise above.

\appendix
\section{Validation, provenance, and rights}
\label{app:provenance}
This paper is a self-contained presentation of the transfer proof,
with published and stipulated inputs identified in
Section~\ref{sec:conventions}. The source revision is
\begin{center}\small\pin\end{center}
of \href{https://github.com/openai/math}{OpenAI Math}.
The accompanying local package records source paths, statement labels,
and hashes, the frozen author packet, its independent proof and
contribution reviews, and both attribution addenda. Reviewer completions
are integrated into the rank lemma, the quantified recovery lemma,
the local-coordinate Jacobian proof, and the literal power-series
hypothesis check. The old exploratory status text in working notes
is not part of the theorem or proof.

The frozen proof packet received the independent verdict
\texttt{PASS\_FULL\_T298P}; the contribution review classified its
mechanism as substantial conditional on correctness and required four
attribution clarifications. Section~\ref{sec:comparison} incorporates
those clarifications. These records concern the frozen analytic proof;
they do not constitute an independent review of this typeset manuscript
or a certification of worldwide novelty. This manuscript is an author
submission for that separate review.

A companion Lean package has four support declarations and four
counterexample controls. It proves a real inner-product coefficient
isometry assuming a column Gram identity, its scaled leading-score
energy identity, a scalar gap rearrangement assuming both analytic
increment bounds and endpoint comparison, and strictness under an
explicit strict error bound. Its controls demonstrate failures with
$k^2$ in place of $k^4$, missing endpoint comparison, a reversed source
gap sign, or a non-strict error threshold. It does not formalize
conjugate variables, analytic inverse series, microstate recovery,
entropy limits, the source theorem, or the theorem of this paper.
No full formalization is claimed. The scalar support statements
quantify over finite real numbers, not extended-real entropy functions.
Their reported axiom dependencies are \texttt{propext},
\texttt{Classical.choice}, and \texttt{Quot.sound}. Reproduction builds
fresh support and control oleans while trusting the installed Lean
toolchain and pre-existing pinned dependency oleans. It rebuilds
neither the compiler nor mathlib; Git source cleanliness does not
authenticate those cached binaries. The independently reviewed
supplement checks ten additional concrete examples and five expected
rejected mutations, without enlarging the analytic coverage.

Finite symbolic and numerical diagnostics in the earlier proof package
check matrix coefficients, the coarse Gram identities, noise covariance,
metric multiplicities, and finite recovery examples. Such diagnostics
are support for calculations, not proofs of the entropy limit theorem.
The local reproducibility package builds this PDF from its complete
LaTeX source, records the compiler and PDF hashes, and supplies a
read-only externally anchored integrity verifier. Integrity checks and
visual inspection verify files and layout, not mathematical truth.
The construction is existential at the fixed large value $k=2^{32}$;
no useful quantitative $\eps$ or low-dimensional numerical witness
has been produced.

The author is Myeongjin Kang, with no institutional affiliation stated.
AI assistance was used in research drafting, review workflows, and
typesetting; the detailed local provenance records the role boundaries.
Original manuscript prose and LaTeX are licensed under
\href{https://creativecommons.org/licenses/by/4.0/}{CC BY 4.0}.
Original reproducibility code is licensed under MIT. Third-party
papers, theorem statements, software, and earlier packet materials
retain their own rights and are credited separately; the licenses here
do not relicense them. No third-party full paper is redistributed in
this manuscript package.

\begin{thebibliography}{99}
\raggedright
\bibitem{source298} OpenAI,
\emph{A finite-entropy separation of microstates and nonmicrostates free
entropy}, September 25, 2026. Pinned OpenAI Math preprint, main theorem,
Sections 2, 4--6.
\href{https://github.com/openai/math/blob/fd4aeeb2ee4fc729c18d98444fed42fd0529eeeb/preprints/A-finite-entropy-separation-of-microstates-and-nonmicrostates-free-entropy-September-25-2026/build/paper.tex}{Immutable source}.
\bibitem{source287} OpenAI,
\emph{An isomorphism of the free group factors}, September 23, 2026.
Pinned OpenAI Math preprint, amplification and entropy-dimension sections.
\href{https://github.com/openai/math/blob/fd4aeeb2ee4fc729c18d98444fed42fd0529eeeb/preprints/An-isomorphism-of-the-free-group-factors-September-23-2026/build/main.tex}{Immutable source}.
\bibitem{generation} OpenAI,
\emph{Relative generation and the generator problem for finite factors},
September 23, 2026. Pinned OpenAI Math preprint, introduction
Theorem \texttt{thm:main}; consequences Corollary
\texttt{cor:entropy-generators}.
\href{https://github.com/openai/math/blob/fd4aeeb2ee4fc729c18d98444fed42fd0529eeeb/preprints/Relative-generation-and-the-generator-problem-for-finite-factors-September-23-2026/build/sections/introduction.tex}{Immutable theorem source}.
\bibitem{Shly} D. Shlyakhtenko,
\emph{A microstates approach to relative free entropy},
\href{https://arxiv.org/abs/math/9812079v1}{arXiv:math/9812079v1} (1998).
Definition 2.3, Proposition 2.8, Theorems 2.19 and 4.1, Proposition 4.2,
and Section 5. Version-specific omissions and typos are discussed in text.
\bibitem{NSSentropy} A. Nica, D. Shlyakhtenko, and R. Speicher,
\emph{Maximality of the microstates free entropy for R-diagonal elements},
\href{https://arxiv.org/abs/math/9809081v1}{arXiv:math/9809081v1} (1998).
Lemmas 2.6, 2.8--2.10 and Theorem 3.1.
\bibitem{NSSfisher} A. Nica, D. Shlyakhtenko, and R. Speicher,
\emph{Some minimization problems for the free analogue of the Fisher
information},
\href{https://arxiv.org/abs/math/9809080v1}{arXiv:math/9809080v1} (1998).
Proposition 4.1.
\bibitem{Vchi} D. Voiculescu,
\emph{The analogues of entropy and of Fisher's information measure in
free probability theory. II}, Invent. Math. \textbf{118} (1994), 411--440.
\href{https://doi.org/10.1007/BF01231539}{doi:10.1007/BF01231539}.
\bibitem{Vstar} D. Voiculescu,
\emph{The analogues of entropy and of Fisher's information measure in
free probability theory. V. Noncommutative Hilbert transforms},
Invent. Math. \textbf{132} (1998), 189--227.
\href{https://doi.org/10.1007/s002220050222}{doi:10.1007/s002220050222}.
\bibitem{Vmatrix} D. Voiculescu,
\emph{Circular and semicircular systems and free product factors},
in \emph{Operator Algebras, Unitary Representations, Enveloping Algebras,
and Invariant Theory}, Progress in Mathematics 92, Birkh\"auser (1990),
45--60. Classical semicircular matrix construction; also recalled in
\cite{NSSentropy,NSSfisher}.
\bibitem{BCG} P. Biane, M. Capitaine, and A. Guionnet,
\emph{Large deviation bounds for matrix Brownian motion},
Invent. Math. \textbf{152} (2003), 433--459.
\href{https://doi.org/10.1007/s00222-002-0281-4}{doi:10.1007/s00222-002-0281-4}.
\bibitem{JP} D. Jekel and J. Pi,
\emph{An elementary proof of the inequality $\chi\leq\chi^*$ for
conditional free entropy}, Doc. Math. \textbf{29} (2024), 1085--1124.
\href{https://doi.org/10.4171/DM/969}{doi:10.4171/DM/969}.
\bibitem{Marinov} M. S. Marinov,
\emph{Invariant volumes of compact groups}, J. Phys. A \textbf{13}
(1980), 3357--3366. The normalization used here is the
$\Rea\Tr(ab^*)$ convention of \cite[Lemma 2.8]{NSSentropy}.
\end{thebibliography}
\end{document}
